% TOP_PROBLEM: {"id": 30006708, "title": "Derandomization of Polynomial Identity Testing", "category_id": 15, "category": {"id": 15, "name": "computer_science", "display_name": "Computer Science", "description": "Computational complexity, algorithms, and theoretical CS.", "slug": "computer-science", "order_index": 15, "created_at": "2026-07-31T15:26:25.671Z"}, "status": "open"}
% FIRST_POSTED: 2026-09-27
% !TeX program = lualatex
% Compile twice: lualatex 30006708.tex
% Standalone research notebook; original section preserved.
% Research additions: 27 September 2026. No external TeX input or BibTeX file.
\documentclass[11pt,a4paper]{article}
\usepackage[margin=24mm,headheight=14pt]{geometry}
\usepackage{fontspec}
\setmainfont{Latin Modern Roman}
\setsansfont{Latin Modern Sans}
\setmonofont{Latin Modern Mono}
\usepackage{amsmath,amssymb,mathtools}
\usepackage{unicode-math}
\setmathfont{Latin Modern Math}
\usepackage{microtype}
\usepackage{enumitem,xurl,needspace}
\usepackage[dvipsnames]{xcolor}
\usepackage{hyperref}
\hypersetup{unicode=true,colorlinks=true,linkcolor=MidnightBlue,urlcolor=MidnightBlue,
 pdftitle={Research notebook - Problem 30006708},pdfauthor={Open Problems Math},bookmarksnumbered=true}
\usepackage{bookmark}
\usepackage{titlesec}
\titleformat{\section}{\Large\bfseries\raggedright}{\thesection}{0.7em}{}
\usepackage{fancyhdr}
\pagestyle{fancy}\fancyhf{}
\fancyhead[L]{\small\sffamily Open Problems Math | Research notebook}
\fancyhead[R]{\small\sffamily Problem 30006708 | 27 September 2026}
\fancyfoot[C]{\thepage}
\setlength{\parindent}{0pt}
\setlength{\parskip}{5pt plus 1pt}
\setlength{\emergencystretch}{3em}
\setcounter{tocdepth}{1}
\setcounter{secnumdepth}{1}
\setlist[itemize]{leftmargin=3em,itemsep=5pt,topsep=4pt,parsep=0pt}
\allowdisplaybreaks
\numberwithin{equation}{section}
\newcommand{\N}{\mathbb{N}}
\newcommand{\Z}{\mathbb{Z}}
\newcommand{\Q}{\mathbb{Q}}
\newcommand{\R}{\mathbb{R}}
\newcommand{\C}{\mathbb{C}}
\newcommand{\F}{\mathbb{F}}
\newcommand{\A}{\mathbb{A}}
\newcommand{\PP}{\mathbb P}
\newcommand{\E}{\mathbb{E}}
\newcommand{\eps}{\varepsilon}
\newcommand{\dd}{\,\mathrm d}
\newcommand{\sref}[2]{\hyperlink{source:#1:#2}{[S#2]}}
\newcommand{\eref}[2]{\hyperlink{extra:#1:#2}{[E#2]}}
\newif\ifshowcatalogue
\showcataloguetrue
\newcommand{\cataloguescope}[1]{\ifshowcatalogue
 \paragraph{Exact scope in the supplied catalogue.}
 \begin{quote}\small #1\end{quote}\fi}
\begin{document}
\setcounter{section}{0}
\section*{\texorpdfstring{Derandomization of Polynomial Identity Testing}{Derandomization of Polynomial Identity Testing}}\label{30006708}
\subsection{Definitions and mathematical statement}
An arithmetic circuit over a specified effective field has variables and field constants as inputs and addition and multiplication gates. Its output is a polynomial $f_C$. Polynomial identity testing asks whether
\[
 f_C\equiv0\quad\text{in the polynomial ring},
\]
not whether it vanishes at one input or has some root. The goal is a deterministic polynomial-time test in the circuit's encoded size and the necessary field-representation parameters. The catalogue leaves the field, characteristic, and constant-encoding convention implicit; these must be fixed before interpreting ``polynomial time.'' General circuits, not just sparse polynomials or formulas of bounded depth, are intended.
\par\smallskip\textit{Formulation and concept references:} \sref{30006708}{1}.
\cataloguescope{Determine whether polynomial identity testing for arithmetic circuits (deciding whether a given arithmetic circuit computes the identically zero polynomial) can be derandomized into deterministic polynomial time.}
\subsection{Short English statement}
Can one deterministically and efficiently recognize when a general arithmetic circuit computes the zero polynomial?
% BEGIN RESEARCH_ADDITION ATTEMPT_70
\subsection{Research attempt}

\paragraph{Current frontier and model convention.}
Kaplan--Shpilka's 2026 read-four result gives deterministic polynomial-time
white-box testing, and quasipolynomial-time black-box testing, for that
restricted formula model in characteristic zero or at least five; it is
not a general-circuit algorithm. \eref{30006708}{1}
The Kabanets--Impagliazzo theorem links an appropriate deterministic
identity test for integer circuits to a disjunction of major circuit lower
bounds. Its field and algorithmic conventions cannot be silently exchanged
with a finite-field statement. \eref{30006708}{2}
Here the concrete attack fixes a finite field $\F_q$ with an effective
representation and unit-cost field arithmetic (or its usual polynomial-bit
implementation). This resolves a local ambiguity for the attempt, not the
unspecified field convention of the whole catalogue target. In particular,
no claim about arbitrary rational constants of enormous intermediate height
is inferred from the calculation.

\paragraph{Attack route.}
Direct grid evaluation costs exponentially many points. Kronecker
substitution compresses many variables into one but produces exponential
degree. The route tested below preserves this compression while evaluating
in many small quotient rings. The hoped-for missing lemma is that small
circuits admit a polynomial-size collection of collision-free residue
certificates even when their expanded supports are large. First a precise
sparse version is proved; then a small circuit is used to test exactly
where that proof ceases to be efficient.

\paragraph{Attempt: compressed exponents with exact collision control.}
Let $C$ have $s$ gates in $n$ variables, and let $D\geq1$ be an upper bound
on every individual degree. A formal degree bound of $2^{O(s)}$ is available
without expanding the circuit. Set $w_i=(D+1)^{i-1}$ and
\[
 F(t)=f_C(t^{w_1},\ldots,t^{w_n}),\qquad M=(D+1)^n.
\]
Base-$(D+1)$ uniqueness shows that distinct exponent vectors with coordinates
at most $D$ have distinct encoded exponents in $[0,M)$. Thus $F=0$ if and
only if $f_C=0$. Although $M$ is large, $\log M=n\log(D+1)$ is polynomial
in the circuit parameters under discussion.

Suppose, as an additional promise, that the nonzero polynomial has at most
$T$ monomials after collecting coefficients. Write
$F(t)=\sum_{j=1}^u c_jt^{e_j}$, with $u\leq T$, $c_j\ne0$ and distinct
$0\leq e_j<M$. Define the nonzero integer
\[
 \Delta=\prod_{i<j}|e_i-e_j|,
 \qquad 1\leq\Delta\leq M^{T(T-1)/2}.
\]
For $u\leq1$ the product is one. A prime $\ell$ at which two exponents
coincide modulo $\ell$ divides $\Delta$. Therefore, among any distinct
primes $\ell_1,\ldots,\ell_R$ satisfying
\begin{equation}\label{eq70:product}
 \prod_{j=1}^R\ell_j>M^{T(T-1)/2},
\end{equation}
at least one separates all the exponents. At that prime,
\begin{equation}\label{eq70:quotient}
 F(t)\not\equiv0\pmod{t^\ell-1}.
\end{equation}
The proof is entirely about integer exponents; it does not assume that
$\ell$ differs from the characteristic of $\F_q$. Even if the quotient ring
has nilpotents, $1,t,\ldots,t^{\ell-1}$ is still an $\F_q$-basis, so distinct
residues with nonzero coefficients cannot cancel.

One can take
$R=1+\lceil T(T-1)\log_2(M)/2\rceil$, since every prime is at least two.
The first $R$ primes have polynomial magnitude in $R$, and can be found by
a sieve; an explicit standard bound is recorded in \eref{30006708}{3}.
For each prime compute $t^{w_i}$ modulo $t^\ell-1$ by reducing $w_i$
modulo $\ell$, then evaluate every gate of $C$ as a vector of $\ell$
coefficients. Addition is componentwise and multiplication is cyclic
convolution, taking $O(\ell^2)$ field operations by the elementary method.
The algorithm accepts identity only if every quotient evaluation vanishes.
An identically zero polynomial certainly passes every test; for a nonzero
polynomial under the support promise, \eqref{eq70:product}--\eqref{eq70:quotient}
reject it. Hence this is a deterministic procedure polynomial in
$s,n,T,\log D$ and the chosen field-operation costs. It need not know the
actual support or the integer $\Delta$.

\paragraph{Attempt: the general-circuit certificate that is still missing.}
A usable strengthening would give a deterministic algorithm that, from $C$,
constructs exponent maps $w^{(j)}$ and moduli $m_j$ such that
\[
 f_C\ne0\ \Longrightarrow\
 f_C(t^{w^{(j)}_1},\ldots,t^{w^{(j)}_n})
          \not\equiv0\pmod{t^{m_j}-1}\quad\text{for some }j,
\]
with the sum of quotient dimensions, all exponent bit lengths, and the
construction time polynomial in the input size. Evaluation of the resulting
certificates would then be polynomial-time by the preceding gate simulation.
This is a sufficient condition, not an asserted equivalent to all possible
white-box derandomizations. In particular, constructing the certificates
cannot itself call an identity oracle or enumerate an exponential support.

The support proof does not give this strengthening. The polynomial
$\prod_{i=1}^n(1+x_i)$ has a circuit of linear size and $2^n$ monomials
over every field. For it the displayed worst-case choice of $R$ is
exponential. The example does not prove that this polynomial needs many
quotients; it proves that the support-size estimate fails to bound the
cost by circuit size. It leaves open whether cancellations and circuit
structure could yield a smaller, differently organized family.

\paragraph{Stress test.}
There is a separate finite-field pitfall: $x^q-x$ is a nonzero polynomial
in $\F_q[x]$ but vanishes at every point of $\F_q$. Testing only scalar
inputs over the base field cannot decide the polynomial question. The
quotient construction avoids that error because it tests coefficients in
a polynomial quotient, not just the induced function on $\F_q$.

A single quotient also fails universally: $t^m-1$ vanishes in the quotient
by $t^m-1$, and small circuits can manufacture very large exponents by
repeated squaring. The collision argument deliberately uses several
moduli. It assumes a valid support bound and a valid degree bound; neither
is inferred from the number of variables alone. Nonassociative-algebra
algorithms in \sref{30006708}{3} use a different algebraic model. They do not
justify forgetting the commutativity and associativity of the target's
polynomial ring.

The companion computation checks several residue collisions and their
removal by another modulus over $\F_5$, including the polynomial/function
distinction. These finite tests audit the mechanism, not the missing
polynomial bound for all circuits. The sparse calculation is elementary
and is not offered as a previously unknown sparse-PIT theorem.

\paragraph{Outcome.}
\textbf{Outcome: Exploratory approach with unresolved gap.}

A complete sparse-support certificate argument and its bit-size accounting
were obtained for the declared finite-field model. The attempt identifies
a concrete generalization that would suffice for deterministic testing,
but an exponential-support circuit breaks the only bound proved here.
No general derandomization or circuit lower bound follows from this notebook.
A continuation must construct polynomial-cost certificates from circuit
structure itself, and must separately address the field and constant
encoding convention of any broader formulation.

% Keep the local bibliography together rather than leave a source fragment.
\needspace{170mm}
% END RESEARCH_ADDITION ATTEMPT_70
\subsection{Sources}
\begin{itemize}\raggedright
\item[\textnormal{[S1]}]\hypertarget{source:30006708:1}{} Derandomizing Polynomial Identity over Finite Fields Implies Super-Polynomial Circuit Lower Bounds for NEXP (Bin Fu).
\url{https://arxiv.org/abs/1311.2358}.
{\small\textit{Catalogue formulation source.}}
\item[\textnormal{[S2]}]\hypertarget{source:30006708:2}{} Polynomial Identity Testing for Read-4 Arithmetic Formulas.
\url{https://drops.dagstuhl.de/storage/00lipics/lipics-vol383-ccc2026/LIPIcs.CCC.2026.25/LIPIcs.CCC.2026.25.pdf}.
{\small\textit{Catalogue research/status reference; not a proof endorsement.}}
\item[\textnormal{[S3]}]\hypertarget{source:30006708:3}{} Efficient Polynomial Identity Testing Over Nonassociative Algebras.
\url{https://arxiv.org/abs/2509.11349}.
{\small\textit{Catalogue research/status reference; not a proof endorsement.}}
% BEGIN RESEARCH_ADDITION REFERENCES_70
\item[\textnormal{[E1]}]\hypertarget{extra:30006708:1}{} Nimrod Kaplan and Amir Shpilka,
\emph{Polynomial Identity Testing for Read-4 Arithmetic Formulas},
ECCC TR26-076 (2026), abstract and model/characteristic hypotheses.
\url{https://eccc.weizmann.ac.il/report/2026/076/}.
{\small\textit{Inspected 27 September 2026; restricted-model result only.}}
\item[\textnormal{[E2]}]\hypertarget{extra:30006708:2}{} Valentine Kabanets and Russell Impagliazzo,
\emph{Derandomizing Polynomial Identity Tests Means Proving Circuit Lower Bounds},
Computational Complexity 13 (2004), 1--46; ECCC TR02-055,
main derandomization-to-lower-bounds theorem.
\url{https://www2.cs.sfu.ca/~kabanets/Research/poly.html}.
{\small\textit{Authors' statement inspected 27 September 2026; the integer-circuit
lower-bound disjunction is not transferred to another field without its hypotheses.}}
\item[\textnormal{[E3]}]\hypertarget{extra:30006708:3}{} J. Barkley Rosser and Lowell Schoenfeld,
\emph{Approximate formulas for some functions of prime numbers},
Illinois Journal of Mathematics 6 (1962), 64--94, bounds for the $r$th prime.
\url{https://projecteuclid.org/journals/illinois-journal-of-mathematics/volume-6/issue-1/Approximate-formulas-for-some-functions-of-prime-numbers/10.1215/ijm/1255631807.full}.
{\small\textit{Only the standard consequence $p_r=O(r\log(2r))$ is used for
sieve cost; not a prime-distribution conjecture.}}
% END RESEARCH_ADDITION REFERENCES_70
\end{itemize}

\end{document}
